\newpage
\section*{Теоретическая часть}
\addcontentsline{toc}{section}{Теоретическая часть}

\subsection*{Постановка задачи}
\addcontentsline{toc}{subsection}{Постановка задачи}

В данной работе точки Лагранжа рассматриваются в рамках ограниченной круговой
задачи трёх тел. Именно эта модель будет основной и в данной работе.
Предполагается, что два тела с массами $m_1$ и $m_2$ движутся
по круговым орбитам вокруг общего центра масс, а третье тело имеет настолько
малую массу, что его влиянием на движение первых двух тел можно пренебречь.

Такое допущение хорошо подходит для описания систем вида Солнце - Земля,
Земля - Луна или планета - космический аппарат. При этом движение малого тела
определяется только гравитацией двух основных тел и выбором системы отсчёта.

Если рассматривать задачу в инерциальной системе, то координаты двух массивных
тел непрерывно меняются со временем. Это делает анализ неудобным. Поэтому
переходят к вращающейся системе координат, связанной с двумя основными телами.
В этой системе они неподвижны, а значит, можно изучать не просто движение
малого тела, но и возможные положения равновесия.

\subsection*{Что такое точки Лагранжа}
\addcontentsline{toc}{subsection}{Что такое точки Лагранжа}

Точки Лагранжа, или точки либрации, - это такие точки во вращающейся системе
отсчёта, в которых равнодействующая всех сил, приложенных к малому телу,
обращается в нуль. Если поместить частицу точно в такую точку и сообщить ей
нулевую относительную скорость, то в рамках идеализированной модели она
останется в этом положении.

Иначе говоря, точки Лагранжа представляют собой стационарные решения
ограниченной круговой задачи трёх тел. Их существование связано с тем, что во
вращающейся системе наряду с гравитационным притяжением учитываются также
инерционные эффекты, возникающие из-за вращения.

Всего существует пять таких точек. Они обозначаются $L_1$, $L_2$, $L_3$,
$L_4$ и $L_5$. Три из них лежат на прямой, проходящей через два основных
тела, а две другие образуют с этими телами равносторонние треугольники.

\subsection*{Уравнения движения во вращающейся системе}
\addcontentsline{toc}{subsection}{Уравнения движения во вращающейся системе}

В безразмерной форме уравнения движения малого тела во вращающейся системе
обычно записываются как
\[
\ddot{x} - 2\dot{y} = \frac{\partial U}{\partial x}, \qquad
\ddot{y} + 2\dot{x} = \frac{\partial U}{\partial y}, \qquad
\ddot{z} = \frac{\partial U}{\partial z},
\]
где $U$ - эффективный потенциал системы. В него входят как гравитационные
члены, так и слагаемое, связанное с вращением системы координат.

Эта функция записывается в виде
\[
U = \frac{x^2 + y^2}{2} + \frac{1-\mu}{r_1} + \frac{\mu}{r_2},
\]
где $\mu$ - безразмерный массовый параметр, а $r_1$ и $r_2$ - расстояния от
малого тела до двух основных тел.

Параметр $\mu$ показывает, как соотносятся массы двух основных тел. Обычно он
вводится как отношение массы меньшего тела к сумме масс обоих тел:
\[
\mu = \frac{m_2}{m_1 + m_2}.
\]
Если одно из тел значительно тяжелее другого, то $\mu$ мало. Например, в
системе Солнце - Юпитер этот параметр очень мал, а в системе Земля - Луна он
заметно больше, но всё равно меньше $\frac{1}{2}$. Именно от значения $\mu$
зависят положение точек либрации и условие устойчивости треугольных точек.

Точки либрации находятся из условия равновесия, когда все правые части
системы обращаются в нуль. Тогда необходимо решить систему
\[
\frac{\partial U}{\partial x} = 0, \qquad
\frac{\partial U}{\partial y} = 0, \qquad
\frac{\partial U}{\partial z} = 0.
\]
Именно из этих условий получаются пять особых точек.

\subsection*{Расположение точек $L_1$-$L_5$}
\addcontentsline{toc}{subsection}{Расположение точек $L_1$--$L_5$}

\textbf{Точка $L_1$.}
Точка $L_1$ расположена между двумя массивными телами. В этой области
гравитационное притяжение обоих тел совместно с вращением системы создаёт
условие равновесия.

\textbf{Точка $L_2$.}
Точка $L_2$ расположена за меньшим из двух массивных тел. Например, в системе
Солнце - Земля эта точка находится за Землёй.

\textbf{Точка $L_3$.}
Точка $L_3$ расположена за более массивным телом, на противоположной стороне
системы. В системе Солнце - Земля она находилась бы по другую сторону Солнца.

\textbf{Точки $L_4$ и $L_5$.}
Точки $L_4$ и $L_5$ расположены так, что вместе с двумя массивными телами
образуют вершины равносторонних треугольников. Одна из них находится впереди
меньшего тела по орбите, а другая - позади него.

\subsection*{Как находятся координаты точек либрации}
\addcontentsline{toc}{subsection}{Как находятся координаты точек либрации}

Из анализа уравнений следует, что для точек либрации
обязательно выполняется условие $z = 0$, то есть все пять точек лежат в
плоскости движения основных тел.

Для треугольных точек решение получается сравнительно просто. Они
соответствуют вершинам правильных треугольников, одна из
сторон которых совпадает с отрезком между двумя основными телами. В
безразмерных координатах они имеют вид
\[
L_4 = \left(\frac{1}{2} - \mu, \frac{\sqrt{3}}{2}, 0\right), \qquad
L_5 = \left(\frac{1}{2} - \mu, -\frac{\sqrt{3}}{2}, 0\right).
\]

Коллинеарные точки $L_1$, $L_2$, $L_3$ находятся сложнее. Для них полагают
$y = 0$ и получают нелинейное уравнение относительно координаты $x$. В явном
простом виде его решения обычно не выписываются, поэтому на практике эти точки
определяют численно. Именно такой подход удобно использовать и в программной
части проекта.

\subsection*{Устойчивость точек}
\addcontentsline{toc}{subsection}{Устойчивость точек}

После нахождения точек либрации естественно возникает вопрос об их
устойчивости. Для этого рассматриваются малые отклонения от положения
равновесия. Пусть точка Лагранжа имеет координаты $(x_0, y_0, 0)$, а реальное
положение малого тела немного отличается от неё. Тогда уравнения движения
линеаризуют в окрестности этой точки и получают систему для малых возмущений.

Смысл такого подхода прост. Если малое отклонение со временем возрастает, то
точка неустойчива. Если же возмущение остаётся ограниченным, то точка может
считаться устойчивой хотя бы в первом приближении.

Точки $L_1$, $L_2$ и $L_3$ являются неустойчивыми. Это означает, что если
поместить малое тело точно в такую точку, то в идеальной модели оно может там
оставаться. Но если немного сдвинуть тело, даже очень слабо, то со временем
оно начнёт удаляться от этой точки.

Для треугольных точек $L_4$ и $L_5$ ситуация иная. Получается условие
устойчивости через массовый параметр $\mu$:
\[
\mu < \mu^* \approx 0.0385.
\]
Если это условие выполнено, то треугольные точки являются устойчивыми в первом
приближении. Это означает, что тело, помещённое в их окрестность с небольшой
относительной скоростью, может достаточно долго оставаться рядом с ними.

Отдельно отметим, что условие устойчивости выполняется, например,
для системы Солнце - Юпитер и для системы Земля - Луна. Поэтому в окрестности
точек $L_4$ и $L_5$ могут существовать естественные объекты. Классическим
примером являются троянские астероиды Юпитера. Для системы Земля - Луна
упоминаются облака Кордылевского вблизи треугольных точек либрации.

Для данного проекта это особенно важно, потому что численное моделирование как
раз и позволит увидеть эту разницу. Если задать начальные условия в окрестности
неустойчивой точки, траектория будет постепенно уходить от неё. Если же задать
начальные условия рядом с устойчивой точкой, движение останется ограниченным в
её окрестности.

\subsection*{Практическое значение точек Лагранжа}
\addcontentsline{toc}{subsection}{Практическое значение точек Лагранжа}

Точки Лагранжа имеют не только теоретический, но и прикладной смысл. Они
используются в небесной механике, космонавтике и задачах проектирования
космических миссий.

Окрестности точек $L_1$ и $L_2$ удобны для размещения научных аппаратов,
поскольку позволяют получать выгодные условия наблюдения. Точки $L_4$ и $L_5$
интересны тем, что в них и рядом с ними могут существовать долгоживущие
конфигурации движения.

\subsection*{Итог}
\addcontentsline{toc}{subsection}{Итог}

Таким образом, ограниченная круговая задача трёх тел приводит к существованию
пяти особых точек равновесия. Три из них являются коллинеарными, две -
треугольными. Эти точки выводятся из уравнений движения во вращающейся системе,
после чего исследуется их устойчивость.

Для дальнейшей вычислительной части особенно важно следующее. Точки $L_1$,
$L_2$, $L_3$ неустойчивы, поэтому малое отклонение от них со временем
нарастает. Точки $L_4$ и $L_5$ при выполнении условия на параметр $\mu$
обладают устойчивостью, поэтому движение в их окрестности может оставаться
ограниченным. Именно это различие и будет проверяться в численных
экспериментах.
